\section{Formulation}\label{sec:formulation}

% Plan and required inputs: notes/outline.md, Section 2. Since session 18 the look-ahead is the
% Taylor polynomial of the flow, as in kinematics_flow_controller.hpp; the staged (symplectic
% Euler) look-ahead of sessions 8-17 survives only in Remark rem:staged. Subsections 2.1, 2.9,
% 2.10 and 2.12 still carry only their plan.

\subsection{Geometric preliminaries}\label{sec:preliminaries}

% Only what the general theory and the three examples need: smooth manifolds, tangent and
% cotangent bundles, Lie derivatives along vector fields (Taylor look-ahead), the Lie groups SO(3)
% and SE(3) with the hat/vee maps (Taylor look-ahead on SO(3), Magnus-type expansion).
\TODO{流形、切丛与余切丛、李导数、李群 SO(3)/SE(3) 及 hat/vee 映射等预备知识，篇幅从简}

\subsection{Problem setting}\label{sec:problem}

We consider plants of the form
\begin{equation}\label{eq:plant}
    \dot\bx = \bv(\bx,\bu), \qquad \bx \in \R^n, \quad \bu \in \mathcal{U} \subseteq \R^m,
\end{equation}
with a smooth vector field $\bv$.
\TODO{本节暂在局部坐标中表述；流形（含 SO(3)）上的内蕴表述待 2.1 节完成后统一}
The state is partitioned into $L$ levels of increasing differential order,
\begin{equation}\label{eq:levels}
    \bx = (\bx_1,\dots,\bx_L), \qquad \bv = (\bv_1,\dots,\bv_L), \qquad \bx_\ell \in \R^{n_\ell},
\end{equation}
so that positions precede velocities, attitudes precede angular velocities, and so on; for $L = 2$
we also write $(\mathbf{q},\mathbf{p}) = (\bx_1,\bx_2)$.
The partition only orders the state components by the number of integrations that separate them
from the input; in particular, $\mathbf{q}$ and $\mathbf{p}$ need not be canonically conjugate.

The reference is given as a sequence of state samples $\bx_{r,0},\bx_{r,1},\dots$, as produced by a
planner, without time stamps.
The evolution times between consecutive samples are roughly stable, close to the update interval
$h$ of the controller, but not exactly uniform, and neither values between the samples nor the
input that generates them are assumed to be available.
At round $k$ of the controller, at time $t_k$, the target is the $m$th next sample,
\begin{equation}\label{eq:xref}
    \bxref(t) = \bx_{r,k+m}, \qquad t\in[t_k,t_{k+1}),
\end{equation}
which the controller aims to reach after the preview time $\tau = mh$, the time in which the
reference is expected to reach it; the integer $m \ge 2$ is a design parameter
(\cref{rem:target}), and taking a sample at a fixed index offset avoids any projection of the state
onto the reference.
The tracking error at round $k$ is the deviation of the state from the sample due at that round,
which is the target set $m$ rounds earlier,
\begin{equation}\label{eq:error}
    \mathbf{e}_k := \bx(t_k) - \bx_{r,k};
\end{equation}
we do not split it into along-track and cross-track components.
For the analysis only, we regard the samples as taken from a trajectory $\bx_r(t)$ of the plant,
$\dot\bx_r = \bv(\bx_r,\bu_r)$ for some input $\bu_r(t)$, with $\bx_{r,k} = \bx_r(t_k)$; the controller
never uses $\bx_r$ or $\bu_r$.
Set-point regulation is the special case of a constant reference.

\subsection{The flow-based cost}\label{sec:cost}

Given a look-ahead map $\Phi_\tau\colon \R^n\times\mathcal{U}\to\R^n$, which predicts the state reached
after the preview time under the input $\bu$, \method{} penalizes the mismatch between the
prediction and the target sample,
\begin{equation}\label{eq:cost}
    F(\bx,\bu,t) = \norm[\big]{\Gamma\bigl(\Phi_\tau(\bx,\bu) - \bxref(t)\bigr)}^2,
\end{equation}
where the invertible matrix $\Gamma$ renders every component dimensionless; we write
$W := \Gamma^\T\Gamma$.
With the mean velocity of the look-ahead, $\bv_\tau(\bx,\bu) := (\Phi_\tau(\bx,\bu) - \bx)/\tau$, the
cost reads
\begin{equation}\label{eq:cost-flow}
    F = \tau^2\,\norm[\Big]{\Gamma\Bigl(\bv_\tau(\bx,\bu) - \frac{\bxref - \bx}{\tau}\Bigr)}^2 ,
\end{equation}
so it compares the flow of the plant over the preview time with the mean velocity that would carry
the current state to the target sample in that time.
This is the counterpart of the cost $(\lambda\bv - \tilde{\bv})^2$ of \citet{lan2004variational}: in
the discretized form of their method, the image of a loop point under the flow over one time step is
compared with the next loop point, the loop tangent $\tilde{\bv}$ being the secant to that point; here
the look-ahead of the plant state is compared with the target sample, and the preview time plays the
role of the time-scaling factor $\lambda$.
We use the forward Euler look-ahead
\begin{equation}\label{eq:fe}
    \Phi_\tau^{\mathrm{FE}}(\bx,\bu) = \bx + \tau\bv(\bx,\bu)
\end{equation}
as a point of comparison below.
In practice $\Gamma$ is diagonal, $\Gamma = \diag(\sqrt{w_1},\dots,\sqrt{w_n})\,\bar\Gamma$, where
$\bar\Gamma$ divides each component by a characteristic scale built from the plant parameters and
$\tau$, and the weights $w_i \ge 0$ are pure numbers.

\subsection{Taylor look-ahead}\label{sec:lookahead}

With the forward Euler look-ahead, an input that drives only the higher levels cannot influence the
predicted lower levels: if $\partial_{\bu}\bv_\ell = 0$, the $\ell$th block row of
$\partial_{\bu}\Phi_\tau^{\mathrm{FE}}$ vanishes, and the error of level $\ell$ drops out of
$\nabla_{\bu}F$.
\method{} therefore predicts the lower levels through the time derivatives of the higher ones.
For a fixed input $\bu$, let $\mathfrak{L}_{\bu}$ denote the Lie derivative along $\bv(\cdot,\bu)$, so
that $\mathfrak{L}_{\bu}^j\bx$ is the $j$th time derivative of the state along the solutions of
\eqref{eq:plant} with the input held at $\bu$; thus $\mathfrak{L}_{\bu}\bx = \bv(\bx,\bu)$ and
$\mathfrak{L}_{\bu}^2\bx = \partial_{\bx}\bv(\bx,\bu)\,\bv(\bx,\bu)$.

\begin{definition}[Taylor look-ahead]\label{def:taylor}
The Taylor look-ahead of order $L$ is
\begin{equation}\label{eq:taylor}
    \hat\Phi_\tau(\bx,\bu) = \sum_{j=0}^{L}\frac{\tau^j}{j!}\,\mathbf{c}_j(\bx,\bu),
\end{equation}
where $\mathbf{c}_j = \mathfrak{L}_{\bu}^j\bx$ for $j \le 2$ and, for $j \ge 3$, $\mathbf{c}_j$ is the
first-order expansion of $\mathfrak{L}_{\bu}^j\bx$ in $\bu$ at $\bu = \mathbf{0}$.
\end{definition}

Up to order two, \eqref{eq:taylor} is the Taylor polynomial of the flow of \eqref{eq:plant} over the
preview time with the input held; from order three on, it keeps only the part of each term that is
affine in the input.
If the input drives only the highest level, $\partial_{\bu}\bv_\ell = 0$ for $\ell < L$, it first
appears in the derivatives of level $\ell$ at order $L - \ell + 1$, so that $L$ is the lowest order at
which it reaches every level.
On an integrator chain, $\hat\Phi_\tau$ is the exact flow.
If the plant is affine in the input, $\bv = \mathbf{f}(\bx) + \sum_i \mathbf{g}_i(\bx)u_i$, and the input
vector fields do not depend on the components they drive, $\partial_{\bx}\mathbf{g}_i\,\mathbf{g}_j =
\mathbf{0}$ for all $i$ and $j$, then $\mathbf{c}_1$ and $\mathbf{c}_2$ are affine in $\bu$ as well, so
that $\hat\Phi_\tau$ is affine in $\bu$ and $F$ is a quadratic function of $\bu$.
The three plants of this paper satisfy this condition; in the quadrotor, for instance, the thrust
acts along the attitude, which no input drives directly.
Expanding \eqref{eq:taylor} in $\tau$ gives
\begin{equation}\label{eq:taylor-expansion}
    \hat\Phi_\tau(\bx,\bu) = \bx + \tau\bv(\bx,\bu) + \tau^2\rho(\bx,\bu,\tau), \qquad
    \rho(\bx,\bu,0) = \tfrac12\,\partial_{\bx}\bv(\bx,\bu)\,\bv(\bx,\bu),
\end{equation}
where $\rho$ is a polynomial in $\tau$; its second-order term is what lets the input reach the levels
below the directly driven ones (\cref{prop:sg}).

\begin{remark}[Staged look-ahead]\label{rem:staged}
A cheaper look-ahead, which needs only the first derivatives of $\bv$, stages a single explicit step
in the order of the symplectic Euler scheme~\citep{hairer2006geometric}: every level is advanced over
the preview time with the already advanced higher levels,
$\bx_\ell^+ = \bx_\ell + \tau\bv_\ell(\bx_1,\dots,\bx_\ell,\bx_{\ell+1}^+,\dots,\bx_L^+,\bu)$, from
$\ell = L$ down to $\ell = 1$.
It also lets the input reach every level, but it holds the advanced higher levels over the whole
preview time.
On an integrator chain it gives $x_\ell^+ = \sum_{j\ge\ell}\tau^{j-\ell}x_j + \tau^{L-\ell+1}u$, with
all coefficients equal to one instead of $1/(j - \ell)!$, so that it overstates the reach of the input
into level $\ell$ by the factor $(L - \ell + 1)!$, and its second-order bias leaves a steady offset
(\cref{ex:offset}).
We use the Taylor look-ahead throughout.
\end{remark}

\subsection{Control law and implementation}\label{sec:extended}

The control input becomes a state of the controller and descends the cost:
\begin{equation}\label{eq:hfc}
    \dot\bx = \bv(\bx,\bu), \qquad \dot\bu = -D(\bx,\bu)\,\nabla_{\bu}F(\bx,\bu,t),
\end{equation}
where $F$ is built on $\hat\Phi_\tau$ and the dimension-balancing coefficients form the diagonal
matrix
\begin{equation}\label{eq:D}
    D = \diag(d_1,\dots,d_m), \qquad
    d_i = \Bigl(2\tau^3\,\norm[\big]{\bar\Gamma\,\partial_{u_i}\bv(\bx,\bu)}^2\Bigr)^{-1}.
\end{equation}
The coefficient $d_i$ is $1/\tau$ times the inverse curvature, with respect to $u_i$, of the
unweighted forward Euler cost $\norm{\bar\Gamma(\bx + \tau\bv - \bxref)}^2$.
It gives $\dot u_i$ the units of a rate of the input, and it scales the gradient so that a residual
on which $u_i$ acts directly, weighted with $w = 1$, decays at the rate $1/\tau$: the infinitesimal
counterpart of a full Newton step taken over the preview time.
The preview time therefore sets the loop gain, and the weights are pure numbers.

The controller runs in discrete time.
At round $k$, it measures the state $\bx(t_k)$ and updates the input by a forward Euler step of
\eqref{eq:hfc}, clamped into $\mathcal{U}$,
\begin{equation}\label{eq:update}
    \bu_{k+1} = \Pi_{\mathcal{U}}\Bigl(\bu_k - h\,D\bigl(\bx(t_k),\bu_k\bigr)\,
        \nabla_{\bu}F\bigl(\bx(t_k),\bu_k,t_k\bigr)\Bigr),
\end{equation}
where $\Pi_{\mathcal{U}}$ clamps every component into its bounds.
The plant moves under the input $\bu_k$, held over $[t_k,t_{k+1})$, and $\bu_{k+1}$ takes over at
$t_{k+1}$, which leaves one round for the computation.
We analyze this sampled-data loop below, and we simulate it with an accurate integration of the
plant.
The update may also be taken from a step of a higher-order scheme applied to \eqref{eq:hfc} from the
measured state, the stage at $t_k + c\,h$ then looking ahead over $\tau - c\,h$; the forward Euler
update is the cheapest, and it is the one we use in the experiments.
\TODO{约束的连续时间表述（投影）}

\subsection{Properties of the flow cost}\label{sec:properties}

\TODO{F 的非负性与零点集}

The flow cost has the structure of a control Hamiltonian, and it generalizes the SG method.
Let $Q(\bx,t) = \frac12\norm{\Gamma(\bx - \bxref(t))}^2$, whose gradient
$\nabla_{\bx}Q = W(\bx - \bxref)$ we denote by $\blambda$, and let
$\omega(\bx,\bu,t) = \partial_t Q(\bx,t) + \nabla_{\bx}Q(\bx,t)^\T\bv(\bx,\bu)$ be the speed of $Q$
along \eqref{eq:plant}, so that $\nabla_{\bu}\omega = (\partial_{\bu}\bv)^\T W(\bx - \bxref)$ is the
direction of the SG law for the goal function $Q$.

\begin{proposition}[Hamiltonian form and relation to the speed-gradient law]\label{prop:sg}
Let $\Gamma$ be fixed, write $K(\bx,\bu) = \frac12\norm{\Gamma\bv(\bx,\bu)}^2$, and define the control
Hamiltonian
\begin{equation}\label{eq:hamiltonian}
    \mathcal{H}(\bx,\bu,\blambda) = \blambda^\T\bv_\tau(\bx,\bu)
    + \frac{\tau}{2}\,\norm[\big]{\Gamma\bv_\tau(\bx,\bu)}^2 .
\end{equation}
\begin{enumerate}
    \item For every look-ahead, $F = 2Q + 2\tau\,\mathcal{H}(\bx,\bu,\nabla_{\bx}Q)$, and therefore
    $\nabla_{\bu}F = 2\tau\,\nabla_{\bu}\mathcal{H}$.
    \item For the forward Euler look-ahead \eqref{eq:fe}, $\mathcal{H} = \omega - \partial_t Q + \tau K$,
    and therefore
    \begin{equation}\label{eq:sg-fe}
        \nabla_{\bu}F = 2\tau\,\nabla_{\bu}\omega + 2\tau^2\,\nabla_{\bu}K.
    \end{equation}
    \item For the Taylor look-ahead \eqref{eq:taylor},
    \begin{equation}\label{eq:sg-taylor}
        \nabla_{\bu}F = 2\tau\,\nabla_{\bu}\omega
        + 2\tau^2\Bigl(\nabla_{\bu}K + \partial_{\bu}\rho(\bx,\bu,0)^\T W(\bx - \bxref)\Bigr)
        + O(\tau^3)
    \end{equation}
    uniformly on compact sets.
    Hence $\dot\bu = -2\tau D\bigl(\nabla_{\bu}\omega + O(\tau)\bigr)$: to first order in $\tau$,
    the \method{} law \eqref{eq:hfc} is the differential SG law with the gain matrix $2\tau D$.
    \item Let $W = \diag(W_1,\dots,W_L)$ be block diagonal, and let level $\ell$ satisfy
    $\partial_{\bu}\bv_\ell = 0$. Then the error $\mathbf{e}_\ell = \bx_\ell - \bx_{\mathrm{ref},\ell}$
    enters neither $\nabla_{\bu}\omega$ nor, for the forward Euler look-ahead, $\nabla_{\bu}F$.
    For the Taylor look-ahead, it enters $\nabla_{\bu}F$, up to terms of order $\tau^3$, through
    \begin{equation}\label{eq:sg-lower}
        \tau^2\Bigl(\sum_{k=1}^{L}\partial_{\bx_k}\bv_\ell\,\partial_{\bu}\bv_k\Bigr)^{\!\T}
        W_\ell\,\mathbf{e}_\ell .
    \end{equation}
\end{enumerate}
\end{proposition}

\begin{proof}
(a) Since $\Phi_\tau - \bxref = (\bx - \bxref) + \tau\bv_\tau$, expanding the square in
\eqref{eq:cost} gives
$F = \norm{\Gamma(\bx - \bxref)}^2 + 2\tau(\bx - \bxref)^\T W\bv_\tau + \tau^2\norm{\Gamma\bv_\tau}^2
= 2Q + 2\tau\,\mathcal{H}(\bx,\bu,\nabla_{\bx}Q)$, and $Q$ does not depend on $\bu$.

(b) For the forward Euler look-ahead, $\bv_\tau = \bv$ and
$(\bx - \bxref)^\T W\bv = \nabla_{\bx}Q^\T\bv = \omega - \partial_t Q$; neither $Q$ nor $\partial_t Q$
depends on $\bu$, which gives \eqref{eq:sg-fe}.

(c) With $\hat\Phi_\tau - \bxref = (\bx - \bxref) + \tau\bv + \tau^2\rho$,
\[
    \nabla_{\bu}F
    = 2\bigl(\tau\,\partial_{\bu}\bv + \tau^2\partial_{\bu}\rho\bigr)^\T W
      \bigl((\bx - \bxref) + \tau\bv + \tau^2\rho\bigr).
\]
The terms of order $\tau$ and $\tau^2$ are
\[
    2\tau(\partial_{\bu}\bv)^\T W(\bx - \bxref) = 2\tau\nabla_{\bu}\omega
    \quad\text{and}\quad
    2\tau^2\bigl((\partial_{\bu}\bv)^\T W\bv + (\partial_{\bu}\rho)^\T W(\bx - \bxref)\bigr),
\]
where $(\partial_{\bu}\bv)^\T W\bv = \nabla_{\bu}K$.
The remaining terms carry at least the factor $\tau^3$.
Since $\rho$ is a polynomial in $\tau$ whose coefficients are continuously differentiable in
$(\bx,\bu)$, replacing $\partial_{\bu}\rho(\bx,\bu,\tau)$ by $\partial_{\bu}\rho(\bx,\bu,0)$ changes the
$\tau^2$ term by $O(\tau^3)$.
All bounds are uniform on compact sets, where the functions involved are continuous.
Multiplying by $D$ gives the form of $\dot\bu$.

(d) Since $W$ is block diagonal, $\nabla_{\bu}\omega = \sum_k (\partial_{\bu}\bv_k)^\T
W_k\mathbf{e}_k$.
If $\partial_{\bu}\bv_\ell = 0$, the term with $k = \ell$ vanishes, so $\mathbf{e}_\ell$ does not appear
in $\nabla_{\bu}\omega$, nor, by \eqref{eq:sg-fe}, in $\nabla_{\bu}F$ for the forward Euler
look-ahead.
For the Taylor look-ahead, $\partial_{\bx_k}\bv_\ell$ does not depend on $\bu$, so that
\eqref{eq:taylor-expansion} gives $\partial_{\bu}\rho_\ell(\bx,\bu,0) =
\frac12\sum_k\partial_{\bx_k}\bv_\ell\,\partial_{\bu}\bv_k$, and the $\ell$th block of the
$\tau^2$ term in \eqref{eq:sg-taylor} is \eqref{eq:sg-lower}.
\end{proof}

\begin{remark}[Hamiltonian and energy-like structure]\label{rem:hamiltonian}
The function \eqref{eq:hamiltonian} has the form of the Hamiltonian of Pontryagin's
principle~\citep{pontryagin1962mathematical}, a running cost plus a costate times the velocity: the
gradient $\blambda = \nabla_{\bx}Q$ of the weighted state error plays the role of the costate, the
mean velocity $\bv_\tau$ of the look-ahead that of the velocity, and $\frac{\tau}{2}\norm{\Gamma
\bv_\tau}^2$, the kinetic energy of the mean velocity in the metric $W$, that of the running cost.
Pontryagin's principle requires an optimal input to minimize the Hamiltonian along the solution of
the adjoint equation.
By \cref{prop:sg}(a), \method{} instead lets the input descend $\mathcal{H}$ continuously, with the
gradient of $Q$ in place of the costate, so that $Q$ plays the role of a value function, and it needs
neither an adjoint equation nor a boundary-value problem; hence the name of the method.
Within a round the target is held, $\partial_t Q = 0$, and $F/2 = Q + \tau\mathcal{H}$ adds $\tau$
times this Hamiltonian to the potential $Q$; for the forward Euler look-ahead, $F/2 = Q + \tau\omega +
\tau^2 K$ is the sum of the potential, its rate along the flow, and a kinetic term.
The input thus descends an energy-like function of the state and of its flow.
The name does not refer to Hamiltonian mechanics: with the state frozen, $F$ decreases along
\eqref{eq:hfc} at the rate $(\nabla_{\bu}F)^\T D\,\nabla_{\bu}F$, so that the closed loop is
dissipative, and no symplectic structure is claimed.
\end{remark}

\begin{remark}\label{rem:sg-levels}
If level $\ell$ is separated from the input by more than one level, the sum in \eqref{eq:sg-lower}
vanishes and $\mathbf{e}_\ell$ enters $\nabla_{\bu}F$ at a higher order in $\tau$; on an integrator
chain, the error of the level that lies $j$ levels below the input enters at order $\tau^{j+1}$
(see the proof of \cref{prop:chain}).
The non-dimensionalization by $\bar\Gamma$, which scales the components of different levels by
different powers of $\tau$, makes these contributions comparable in size.
\end{remark}

\subsection{Convergence analysis}\label{sec:convergence}

\TODO{一般非线性系统的收敛性主定理，研究方案见 notes/convergence.md}

\paragraph{Integrator chains.}
The integrator chain is the prototype of the level structure \eqref{eq:levels}; on it, the Taylor
look-ahead is the exact flow, and the \method{} loop can be analyzed exactly.

\begin{proposition}[\method{} on an integrator chain]\label{prop:chain}
Consider the chain $\dot x_1 = x_2,\dots,\dot x_{L-1} = x_L$, $\dot x_L = u$ with scalar states and
input, the Taylor look-ahead \eqref{eq:taylor}, the coefficient \eqref{eq:D}, and
\[
    \Gamma = \ell_0^{-1}\diag\bigl(\sqrt{w_1},\ \sqrt{w_2}\,\tau,\ \dots,\ \sqrt{w_L}\,\tau^{L-1}\bigr)
\]
with a characteristic length $\ell_0$ and weights $w_k > 0$.
For set-point regulation, $\bx_r = 0$ and $u_r = 0$, the extended system \eqref{eq:hfc} is linear
and its characteristic polynomial is $\tau^{-(L+1)}P(\tau s)$, where
\begin{equation}\label{eq:chain-poly}
    P(\sigma) = \sigma^{L+1} + \sum_{k=1}^{L}\frac{w_k}{(L - k + 1)!}\,\sigma^{k-1}e_{L-k+1}(\sigma),
    \qquad e_n(\sigma) = \sum_{i=0}^{n}\frac{\sigma^i}{i!} .
\end{equation}
Consequently, the closed loop is exponentially stable if and only if $P$ is Hurwitz, a condition on
the weights alone, and every closed-loop eigenvalue is $1/\tau$ times a root of $P$.
\end{proposition}

\begin{proof}
Since the chain is linear, \eqref{eq:taylor} is its exact flow,
$x_k^+ = \sum_{j=k}^{L}\tau^{j-k}x_j/(j-k)! + \tau^{L-k+1}u/(L-k+1)!$.
In the scaled variables $z_j = \tau^{j-1}x_j$ and $\nu = \tau^L u$ this reads $\tau^{k-1}x_k^+ = Y_k$
with $Y_k = \sum_{j\ge k} z_j/(j-k)! + \nu/(L-k+1)!$, so that $F = \ell_0^{-2}\sum_{k} w_k Y_k^2$.
Here $\bar\Gamma\,\partial_u\bv = \tau^{L-1}/\ell_0$, so \eqref{eq:D} gives
$d = \ell_0^2/(2\tau^{2L+1})$ and
$\dot u = -d\,\partial_u F = -\tau^{-(L+1)}\sum_k w_k Y_k/(L-k+1)!$.
In the time $\sigma = t/\tau$, with $(\cdot)' = \dd(\cdot)/\dd\sigma$, the chain becomes
$z_j' = z_{j+1}$ for $j < L$ and $z_L' = \nu$, and the control equation becomes
$\nu' = -\sum_k w_k Y_k/(L-k+1)!$.
Since $z_j = z_1^{(j-1)}$ and $\nu = z_1^{(L)}$, the Laplace transform turns $Y_k$ into
$s^{k-1}e_{L-k+1}(s)$ times the transform of $z_1$, and the characteristic polynomial of the
resulting scalar equation for $z_1$ is $P$.
The change of variables is invertible for $\tau > 0$, and the time scaling multiplies the
eigenvalues by $1/\tau$.
\end{proof}

\begin{remark}[Weight design and update interval]\label{rem:chain}
For $L = 1$, $P = \sigma^2 + w_1\sigma + w_1$ is Hurwitz for every $w_1 > 0$.
For $L = 2$,
$P(\sigma) = \sigma^3 + \bigl(\frac{w_1}{4} + w_2\bigr)\sigma^2 + \bigl(\frac{w_1}{2} + w_2\bigr)\sigma
+ \frac{w_1}{2}$, which by the Routh--Hurwitz criterion is Hurwitz if and only if
$\bigl(\frac{w_1}{4} + w_2\bigr)\bigl(\frac{w_1}{2} + w_2\bigr) > \frac{w_1}{2}$.
Equal weights are admissible up to $L = 4$, which covers the chain from the rotor torques to the
position of a quadrotor, but not for $L = 5$, where $P$ with $w_k = 1$ has roots in the right
half-plane; longer chains call for a deliberate choice of the weights.
Since the closed-loop eigenvalues scale with $1/\tau$, the sampled-data loop \eqref{eq:update} is
stable only for $h/\tau$ below a threshold set by the roots of $P$.
For $L = 2$ with unit weights, it is stable if and only if $\tau > 1.526\,h$; a simulation that
advances the plant by the same forward Euler step would show the threshold $\tau > 1.408\,h$ instead.
Both are met by $\tau = mh$ with $m \ge 2$.
\end{remark}

\paragraph{Tracking error caused by the look-ahead.}
Even when the plant is on the reference, the look-ahead generally misses the target sample.
We call
\begin{equation}\label{eq:defect}
    \bdelta(t) := \bxref(t) - \hat\Phi_\tau\bigl(\bx_r(t),\bu_r(t)\bigr)
\end{equation}
the look-ahead defect: the amount by which the look-ahead misses the target when the plant is on the
reference and applies the reference input.
The defect is an object of the analysis only; the controller never evaluates it.

\begin{lemma}[Look-ahead defect]\label{lem:defect}
Let the reference stay in a compact set, let $\bx_r$ be three times and $\bu_r$ once continuously
differentiable, with bounded derivatives. If the target sample is the reference state at the time
$t + \theta$, $\bxref(t) = \bx_r(t+\theta)$ with $0 < \theta \le \bar\theta$, then, uniformly in $t$,
\begin{equation}\label{eq:defect-expansion}
    \bdelta(t) = (\theta - \tau)\,\dot\bx_r(t)
    + \tfrac12\bigl(\theta^2 - \tau^2\bigr)\ddot\bx_r(t)
    + \tfrac12\tau^2\,\partial_{\bu}\bv\bigl(\bx_r(t),\bu_r(t)\bigr)\,\dot\bu_r(t)
    + O(\theta^3 + \tau^3).
\end{equation}
For $\theta = \tau$, the defect is of second order only through the rate of the reference input, which
acts on the directly driven levels, and of third order otherwise.
\end{lemma}

\begin{proof}
Taylor expansion gives
$\bx_r(t+\theta) = \bx_r + \theta\dot\bx_r + \frac12\theta^2\ddot\bx_r + O(\theta^3)$, and
\eqref{eq:taylor-expansion} gives
$\hat\Phi_\tau(\bx_r,\bu_r) = \bx_r + \tau\bv + \frac12\tau^2\partial_{\bx}\bv\,\bv + O(\tau^3)$, with
$\bv$ and its derivatives evaluated at $(\bx_r,\bu_r)$.
Along the reference, $\dot\bx_r = \bv$ and $\ddot\bx_r = \partial_{\bx}\bv\,\bv + \partial_{\bu}\bv\,
\dot\bu_r$; subtracting yields \eqref{eq:defect-expansion}.
\end{proof}

The first term of \eqref{eq:defect-expansion} is of first order: unless the preview time matches the
time $\theta$ at which the reference reaches the target sample, the defect is dominated by a
mismatch of timing, which the closed loop absorbs as a time delay (\cref{ex:offset}).
With the next sample as target, $\theta$ is the sample interval, close to $h$, whereas the update
\eqref{eq:update} requires $\tau > 1.526\,h$ for a double integrator (\cref{rem:chain}).

A defect acts on the closed loop as a persistent disturbance, and so does the rate of the reference
input, which the control state has to follow.

\begin{theorem}[Ultimate bound on the tracking error]\label{thm:offset}
Let $\bxref = \hat\Phi_\tau(\bx_r,\bu_r) + \bdelta$ be a target with defect $\bdelta$, let
$\bxi = (\bx - \bx_r,\ \bu - \bu_r)$, and call nominal the error system obtained from \eqref{eq:hfc}
with $\bdelta = 0$ and without the term $-\dot\bu_r$; $\bxi = 0$ is an equilibrium of the nominal
system.
Suppose that this equilibrium is uniformly exponentially stable on a ball $B_r$, with a Lyapunov
function $V(t,\bxi)$ that satisfies $c_1\norm{\bxi}^2 \le V \le c_2\norm{\bxi}^2$,
$\dot V \le -c_3\norm{\bxi}^2$ along the nominal system, and $\norm{\partial_{\bxi}V} \le
c_4\norm{\bxi}$ on $B_r$. Let
\begin{equation}\label{eq:gbar}
    \bar g := \sup_{t,\ \bxi\in B_r}
        \norm[\big]{2D\,(\partial_{\bu}\hat\Phi_\tau)^\T W\bdelta(t)}
        + \sup_{t}\norm{\dot\bu_r(t)} .
\end{equation}
If $\bar g < \vartheta\,(c_3/c_4)\sqrt{c_1/c_2}\,r$ for some $\vartheta\in(0,1)$, then every solution of
\eqref{eq:hfc} with $\norm{\bxi(t_0)} < \sqrt{c_1/c_2}\,r$ satisfies, after a finite time,
\begin{equation}\label{eq:offset-bound}
    \norm{\bxi(t)} \le \frac{c_4}{c_3}\sqrt{\frac{c_2}{c_1}}\;\frac{\bar g}{\vartheta}.
\end{equation}
\end{theorem}

\begin{proof}
Since $\nabla_{\bu}F = 2(\partial_{\bu}\hat\Phi_\tau)^\T W\bigl(\hat\Phi_\tau(\bx,\bu) -
\hat\Phi_\tau(\bx_r,\bu_r) - \bdelta\bigr)$, the error dynamics are the nominal system plus the
perturbation
$\mathbf{g} = \bigl(\mathbf{0},\ 2D(\partial_{\bu}\hat\Phi_\tau)^\T W\bdelta - \dot\bu_r\bigr)$,
whose norm is bounded by $\bar g$ on $B_r$.
The claim is the ultimate-boundedness result for nonvanishing perturbations of exponentially stable
systems~\citep[Sect.~9.2]{khalil2002nonlinear}.
\end{proof}

By \cref{lem:defect}, the look-ahead leaves a tracking error driven by a disturbance of order
$\abs{\theta - \tau} + \tau^2\sup_t\norm{\dot\bu_r} + \tau^3$.
The constants $c_1,\dots,c_4$ and the coefficients $D$ depend on $\tau$, however, and the order of the
resulting offset depends on how they scale; for the integrator chain it can be computed exactly.
\TODO{一般系统中偏差阶数的严格证明（在 σ = t/τ 归一化时间下处理常数），见 notes/convergence.md}

\begin{example}[Delay of the double integrator]\label{ex:offset}
Take $L = 2$ in \cref{prop:chain} and a reference with constant acceleration $u_r = a_r$, whose
target sample at time $t$ is $\bx_r(t+\theta)$.
Substituting $x_1(t) = x_{1,r}(t - d) + c$, $x_2(t) = x_{2,r}(t - d)$ and $u = a_r$ into the
residuals of the cost shows that both vanish for all $t$ if and only if $d = \tau - \theta$ and
$c = 0$, whatever the weights, because the look-ahead is exact for a constant acceleration.
In steady state the plant therefore follows the reference with the time delay $\tau - \theta$ and
without offset; for $\theta = \tau$ it tracks the reference exactly.
The staged look-ahead of \cref{rem:staged} would leave the offset $-\frac12 a_r\tau^2$ instead.
\end{example}

\begin{remark}[Choice of the target sample]\label{rem:target}
With the next sample as target, $m = 1$, the time to the target is $\theta \approx h$, whereas the
update \eqref{eq:update} requires $\tau > 1.526\,h$ for a double integrator (\cref{rem:chain}); the
error \eqref{eq:error} is then dominated by the delay $\tau - h$, about $(\tau - h)\abs{\dot\bx_r}$ in
the units of the state.
Taking the $m$th next sample with $\tau = mh$, $m \ge 2$, matches the preview time to the time to the
target, so that the first-order term of \eqref{eq:defect-expansion} reduces to the deviation of the
sum of the next $m$ sample intervals from $mh$, while the update interval remains admissible
(\cref{rem:chain}).
What remains are this jitter of the sample intervals, the rate of the reference input, and terms of
third order.
\end{remark}

\subsection{Non-uniform sample intervals}\label{sec:nonuniform}

In the setting of \cref{sec:problem}, the target sample becomes due $m$ rounds after it is set, so
that $\tau = mh$ is the time the controller has to reach it, whatever the evolution times of the
reference between the samples.
If these times are markedly non-uniform, consuming one sample per round makes the samples
inconsistent with the clock of the controller: to reach every sample on time, the plant would have
to move faster or slower than the samples prescribe, and no choice of the preview time removes the
resulting error \eqref{eq:error}.
When the samples represent a trajectory $\bx_r(t)$ that is to be followed in time, their evolution
times can instead be estimated from the samples themselves.
If the rate of the lowest level does not depend on the input and stays away from zero, the time
between two samples is, up to terms of third order in the interval,
\begin{equation}\label{eq:interval}
    \widehat{\Delta t}_k = \frac{2\,\norm{\bx_{1,r,k+1} - \bx_{1,r,k}}}
        {\norm{\bv_1(\bx_{r,k})} + \norm{\bv_1(\bx_{r,k+1})}},
\end{equation}
the distance between the samples over their mean speed.
Summing these estimates gives the times $\hat t_j$ of the samples; at round $k$, the target is then the
first sample with $\hat t_j \ge t_k + mh$, the preview time is $\tau = \hat t_j - t_k$, and the
tracking error is the deviation from $\bx_r(t_k)$.
This selection needs neither time stamps nor a projection of the state onto the reference, and the
coefficients \eqref{eq:D} follow the preview time from round to round.
\TODO{仿真：按估计时间选目标时，相对 $\bx_r(t)$ 的误差与间隔抖动无关（notes Q39），数值在 Experiments 中报告}

\subsection{Design procedure}\label{sec:procedure}

\TODO{设计步骤（Algorithm 框）与参数选择：Taylor 预瞄阶数 L、$\tau = mh$（$m \ge 2$，满足\cref{rem:chain}的 τ > 1.526h）、权重（\cref{rem:chain}的 Hurwitz 条件）、前向 Euler 更新；自行车另见\cref{rem:bicycle-speed}}

\subsection{Example: swing-up and stabilization of a double inverted pendulum on a cart}
\label{sec:example-dipc}

\TODO{动力学模型（输入为作用在小车上的力）；两级层级，Taylor 预瞄阶数为 2；单一控制律完成从下垂位置起摆到收敛于倒立平衡点；控制律推导}

\subsection{Example: trajectory tracking with a kinematic bicycle model}
\label{sec:example-bicycle}

% Session 18: Taylor look-ahead, as in kinematics_flow_controller.hpp (refactored 2026-10-26);
% look-ahead, columns and rates checked against the code (notes/discussion.md, Q40).

\paragraph{Model.}
We describe the vehicle by the kinematic bicycle model~\citep{rajamani2012vehicle,kong2015kinematic},
with the reference point on the longitudinal axis at the distance $l_r$ ahead of the rear axle and
$l_f$ behind the front axle, for instance at the center of gravity; $l = l_f + l_r$ is the
wheelbase.
The state $\bx = (x,y,\psi,\delta,u)$ collects the position of the reference point in the world
frame, the heading, the steering angle of the front wheel, and the longitudinal velocity; the input
$\bu = (a,\omega)$ collects the longitudinal acceleration and the steering rate, which are bounded,
$\mathcal{U} = [a_{\min},a_{\max}]\times[\omega_{\min},\omega_{\max}]$; in this subsection, $\omega$
is the steering rate, not the speed of $Q$ in \cref{sec:properties}.
If neither wheel slips sideways, the yaw rate and the lateral velocity of the reference point follow
from $u$ and $\delta$ as $r = u\tan\delta/l$ and $v = l_r r$, and
\begin{equation}\label{eq:bicycle}
    \dot x = u\cos\psi - v\sin\psi, \quad
    \dot y = u\sin\psi + v\cos\psi, \quad
    \dot\psi = r, \quad
    \dot\delta = \omega, \quad
    \dot u = a,
\end{equation}
with $\abs{\delta} < \pi/2$.
The model is accurate as long as the lateral acceleration stays moderate, so that the tires barely
slip~\citep{kong2015kinematic,polack2017kinematic}.
Because the steering angle and the speed are states, the inputs are their rates, and the bounds of
$\mathcal{U}$ are the rate limits of the actuators.
The state falls into three levels, $\bx_1 = (x,y)$, $\bx_2 = \psi$ and $\bx_3 = (\delta,u)$: the
input drives the third level, the third level drives the heading, and the position is driven by the
heading and, through $u$ and $v$, directly by the third level.
The heading is a continuous function of time, for the plant as for the samples, so that heading
differences need no wrapping.

\paragraph{Look-ahead and cost.}
The Taylor look-ahead \eqref{eq:taylor} of \eqref{eq:bicycle} has order three.
Let $\mathbf{p} = (x,y)$, let $R(\psi)$ be the rotation by the heading and $J = R(\pi/2)$, and write
$\mathbf{b} = (u,v)$ for the velocity of the reference point in the body frame,
$\dot r = (a\tan\delta + u\,\omega\sec^2\delta)/l$ for the yaw acceleration, and
$\dot{\mathbf{b}} = (a, l_r\dot r)$ for the rate of $\mathbf{b}$.
Then
\begin{equation}\label{eq:bicycle-lookahead}
\begin{aligned}
    \delta^+ &= \delta + \tau\omega, \qquad u^+ = u + \tau a, \qquad
    \psi^+ = \psi + \tau r + \tfrac12\tau^2\dot r, \\
    \mathbf{p}^+ &= \mathbf{p} + R(\psi)\Bigl(\tau\mathbf{b}
    + \tfrac12\tau^2\bigl(\dot{\mathbf{b}} + rJ\mathbf{b}\bigr)
    + \tfrac16\tau^3\bigl(2rJ\dot{\mathbf{b}} + \dot r J\mathbf{b} - r^2\mathbf{b}\bigr)\Bigr).
\end{aligned}
\end{equation}
The terms up to second order are the derivatives of the flow, the velocity of the reference point
turning with the heading; at third order, the second derivative $\ddot r$ of the yaw rate, which is
quadratic in $(a,\omega)$, is omitted in the position and in the heading.
With the steering angle held, $\omega = 0$, the look-ahead is therefore exact through $\tau^3$, and
along a straight line under constant acceleration it is exact.
Every component of \eqref{eq:bicycle-lookahead} is affine in $(a,\omega)$.

The target $\bxref = (x_{\mathrm{ref}},y_{\mathrm{ref}},\psi_{\mathrm{ref}},\delta_{\mathrm{ref}},
u_{\mathrm{ref}})$ is the $m$th next sample \eqref{eq:xref}, with $\tau = mh$, and the residuals are
weighted as
\begin{equation}\label{eq:bicycle-cost}
    F = w_p\,\frac{(x^+ - x_{\mathrm{ref}})^2 + (y^+ - y_{\mathrm{ref}})^2}{l^2}
    + w_\psi\,(\psi^+ - \psi_{\mathrm{ref}})^2 + w_\delta\,(\delta^+ - \delta_{\mathrm{ref}})^2
    + w_u\,\frac{\tau^2(u^+ - u_{\mathrm{ref}})^2}{l^2},
\end{equation}
that is, $\bar\Gamma = \diag(l^{-1},l^{-1},1,1,\tau/l)$ in \eqref{eq:cost}: positions are measured
in wheelbases, and the speed by the distance it covers in the preview time, again in wheelbases.
Equal weights on $x$ and $y$ keep $F$ invariant under rotations of the world frame.
Since $\partial_a\bv$ and $\partial_\omega\bv$ are the unit vectors of the components $u$ and
$\delta$, the coefficients \eqref{eq:D} are $d_a = l^2/(2\tau^5)$ and $d_\omega = 1/(2\tau^3)$.

Because the look-ahead is affine in the input,
$\nabla_{\bu}F = 2(\partial_{\bu}\hat\Phi_\tau)^\T W(\hat\Phi_\tau - \bxref)$ with
$W = \diag(w_p/l^2, w_p/l^2, w_\psi, w_\delta, w_u\tau^2/l^2)$, and the columns of
$\partial_{\bu}\hat\Phi_\tau$, with the components ordered as $(x,y,\psi,\delta,u)$, are
\begin{equation}\label{eq:bicycle-columns}
\begin{aligned}
    \partial_a\hat\Phi_\tau &= \tau^2\Bigl(
        \frac12 R(\psi)\begin{bmatrix} 1 - \tau v\tan\delta/l \\ (l_r + \tau u)\tan\delta/l \end{bmatrix},\
        \frac{\tan\delta}{2l},\ 0,\ \frac{1}{\tau}\Bigr), \\
    \partial_\omega\hat\Phi_\tau &= \tau^2\Bigl(
        \frac{u\sec^2\delta}{l}\,R(\psi)\begin{bmatrix} -\tau v/2 \\ l_r/2 + \tau u/6 \end{bmatrix},\
        \frac{u\sec^2\delta}{2l},\ \frac{1}{\tau},\ 0\Bigr).
\end{aligned}
\end{equation}
The \method{} law \eqref{eq:hfc} becomes
\begin{equation}\label{eq:bicycle-law}
    \dot a = -\frac{l^2}{\tau^5}\,(\partial_a\hat\Phi_\tau)^\T W\bigl(\hat\Phi_\tau - \bxref\bigr),
    \qquad
    \dot\omega = -\frac{1}{\tau^3}\,(\partial_\omega\hat\Phi_\tau)^\T W\bigl(\hat\Phi_\tau - \bxref\bigr),
\end{equation}
which the controller integrates by the update \eqref{eq:update}.
A round of the controller thus takes a fixed number of arithmetic operations and a few
trigonometric functions; it involves neither a matrix inversion nor an iteration.

\paragraph{How the errors reach the input.}
With the forward Euler look-ahead, only the residuals of $\delta$ and $u$ would enter
$\nabla_{\bu}F$, and the controller would follow the steering angle and the speed of the samples
without any feedback from heading and position (\cref{prop:sg}).
By \eqref{eq:bicycle-columns}, the Taylor look-ahead lets the heading error enter at order $\tau^2$,
through the dependence of the yaw rate on $\delta$ and $u$, and the position error as well, through
the dependence of the velocity of the reference point on $u$ and, sideways, on $\delta$.
The lateral entry of $\partial_\omega\hat\Phi_\tau$ in the body frame,
$\tau^2 u\sec^2\delta\,(l_r/2 + \tau u/6)/l$, combines two effects: the lateral velocity of the
reference point, which builds up over the preview time with the lever arm $l_r/2$, and the rotation of
the velocity by the heading change, with the lever arm $\tau u/6$.
The second effect, of third order, dominates whenever $\tau u > 3l_r$, and it is the only one when the
reference point is at the rear axle, $l_r = 0$; a look-ahead truncated at second order would leave
the lateral position of such a vehicle without any influence on the steering rate.

\paragraph{Straight references.}
About a straight reference traversed at constant speed, the longitudinal motion reduces to the
integrator chain of \cref{prop:chain}, and the lateral motion to a loop of fourth order.

\begin{proposition}[Bicycle on a straight reference]\label{prop:bicycle}
Linearize the closed loop \eqref{eq:bicycle}, \eqref{eq:bicycle-law} about a straight reference
traversed at the constant speed $U > 0$, whose target at the time $t$ is its state at $t + \tau$, and
let $\kappa = U\tau/l$ be the preview distance in wheelbases.
\begin{enumerate}
    \item The linearization splits into a longitudinal part, in the along-track position, $u$ and $a$,
    and a lateral part, in the cross-track position, $\psi$, $\delta$ and $\omega$.
    \item The longitudinal part is the loop of \cref{prop:chain} with $L = 2$, $\ell_0 = l$,
    $w_1 = w_p$ and $w_2 = w_u$.
    \item The characteristic polynomial of the lateral part is $\tau^{-4}P(\tau s)$ with
    \begin{equation}\label{eq:bicycle-lateral}
        P(\sigma) = \sigma^4 + a_3\sigma^3 + a_2\sigma^2 + a_1\sigma + a_0,
    \end{equation}
    where, with $\beta = \kappa\,(\kappa/6 + l_r/(2l))$,
    \[
    \begin{aligned}
        a_0 &= w_p\beta\kappa^2, &
        a_1 &= w_p\beta\kappa\,(\kappa + l_r/l) + \tfrac12 w_\psi\kappa^2, \\
        a_2 &= w_p\beta\kappa\,(\kappa/2 + l_r/l) + \tfrac12 w_\psi\kappa^2 + w_\delta, &
        a_3 &= w_p\beta^2 + \tfrac14 w_\psi\kappa^2 + w_\delta .
    \end{aligned}
    \]
    \item For unit weights, both parts are exponentially stable for every $\kappa > 0$ and every
    $l_r\in[0,l]$, so that the reference is locally exponentially stable.
\end{enumerate}
\end{proposition}

\begin{proof}
We place the reference on the $x$-axis, where $\psi = \delta = v = 0$, $u = U$ and $a = \omega = 0$,
and denote the deviations from it by the same letters.
The reference input lies in the interior of $\mathcal{U}$, so the clamping is inactive near it.

(a) The lateral deviations enter the longitudinal components of \eqref{eq:bicycle} and
\eqref{eq:bicycle-lookahead} only in products of deviations, and vice versa.
In \eqref{eq:bicycle-columns}, the entries that couple the two parts carry the factor $\tan\delta$
or $v$, so they are of first order, and their products with the residuals drop out of the
linearization of \eqref{eq:bicycle-law}.

(b) The longitudinal part reads $\dot x = u$, $\dot u = a$, $x^+ = x + \tau u + \frac12\tau^2 a$ and
$u^+ = u + \tau a$, with the cost $w_p(x^+ - x_{\mathrm{ref}})^2/l^2 +
w_u\tau^2(u^+ - u_{\mathrm{ref}})^2/l^2$ and $d_a = l^2/(2\tau^5)$.
In the deviations from the reference, whose target is exactly its own look-ahead, this is the loop
of \cref{prop:chain}.

(c) Let $\eta = y/l$, $\nu = \tau\omega$, and $(\cdot)' = \dd(\cdot)/\dd\sigma$ with $\sigma = t/\tau$.
To first order, $\dot r = U\omega/l$ and $v = l_r U\delta/l$, so the lateral part reads
$\eta' = \kappa(\psi + (l_r/l)\delta)$, $\psi' = \kappa\delta$ and $\delta' = \nu$, and its look-ahead
\[
    \frac{y^+}{l} = \eta + \kappa\psi + \kappa\Bigl(\frac{\kappa}{2} + \frac{l_r}{l}\Bigr)\delta
    + \beta\nu,
    \qquad \psi^+ = \psi + \kappa\delta + \tfrac12\kappa\nu, \qquad \delta^+ = \delta + \nu .
\]
Since $\partial_\omega = \tau\partial_\nu$, the control equation is
$\nu' = \tau^2\dot\omega = -\frac12\partial_\nu F
= -w_p\beta\,y^+/l - \frac12 w_\psi\kappa\psi^+ - w_\delta\delta^+$.
In the Laplace domain, $\nu = s\delta$, $\psi = \kappa\delta/s$ and
$\eta = \kappa(\kappa + (l_r/l)s)\delta/s^2$; substituting these into the control equation and
multiplying by $s^2$ gives $P(s)\delta = 0$.
The scaling of time multiplies the eigenvalues by $1/\tau$.

(d) For unit weights, the longitudinal part satisfies
$\bigl(\frac{w_1}{4} + w_2\bigr)\bigl(\frac{w_1}{2} + w_2\bigr) = \frac{15}{8} > \frac{w_1}{2}$
(\cref{rem:chain}).
For the lateral part, the coefficients $a_0,\dots,a_3$ are positive, and the Hurwitz determinants
$a_3a_2 - a_1$ and $a_1(a_3a_2 - a_1) - a_3^2a_0$ are polynomials in $\kappa$ and $l_r/l$;
quantifier elimination over the reals by cylindrical algebraic decomposition~\citep{collins1975quantifier}
shows that both are positive for all $\kappa > 0$ and $l_r/l\in[0,1]$.
Hence $P$ is Hurwitz, and Lyapunov's indirect method gives the local exponential stability.
\end{proof}

\begin{remark}[Speed range of the sampled-data loop]\label{rem:bicycle-speed}
With $\tau = mh$ and $m \ge 2$, the longitudinal part meets the condition $\tau > 1.526\,h$ of
\cref{rem:chain}.
The lateral part bounds the preview distance instead.
For unit weights and $l_r = l/2$, the sampled-data loop \eqref{eq:update}, linearized about a straight
reference, has spectral radius below one for $\kappa < 0.92$, $1.21$ and $3.37$ with $m = 2$, $3$ and
$4$; for $l_r = 0$, the bounds are $1.15$, $1.49$ and $3.94$.
These bounds depend on $\kappa$, $l_r/l$ and $m$ only; the loop can regain stability beyond them, but
we do not rely on it.
In practice, the preview distance $U\tau$ should therefore stay below about one wheelbase; at higher
speeds, the update interval $h$ has to be shortened or the weight $w_p$ lowered.
\end{remark}

\paragraph{Tracking error.}
By \cref{lem:defect} with $\theta = \tau$, the defect is
$\frac12\tau^2\partial_{\bu}\bv\,\dot\bu_r + O(\tau^3)$: it is of second order only in the steering
angle and the speed, $\frac12\tau^2\dot\omega_r$ and $\frac12\tau^2\dot a_r$, and of third order
otherwise.
On references with a held steering angle and a constant acceleration, which include straight lines and
circles traversed at constant speed, the look-ahead is exact through $\tau^3$ and the defect is of
fourth order, so that the steady offsets essentially vanish, whereas the staged look-ahead of
\cref{rem:staged} leaves the vehicle about $\frac12\tau^2$ times the centripetal acceleration outside
a curve.
With $\tau = mh$, the error \eqref{eq:error} is then dominated by the jitter of the sample intervals,
which enters at first order (\cref{rem:target}).
\TODO{稳态偏差的仿真数值（圆周、匀加速直线）在 Experiments 中报告}

\paragraph{Conditions on the reference.}
The samples carry all five components of the state; the controller uses no input of the reference.
For the analysis, the samples are taken from a solution of \eqref{eq:bicycle} with continuously
differentiable inputs in the interior of $\mathcal{U}$, $\abs{\delta_r} < \pi/2$, and a speed bounded
away from zero.
The last condition is essential: as $U \to 0$, $\kappa \to 0$ and the lateral feedback in
\eqref{eq:bicycle-lateral} vanishes with $a_0$ and $a_1$.
This is not a shortcoming of the particular law: by Brockett's necessary
condition~\citep{brockett1983asymptotic}, no continuous time-invariant feedback, static or dynamic,
asymptotically stabilizes a parked pose of the vehicle, since no state near rest and no input
produce a sideways velocity.

\subsection{Example: trajectory tracking of a quadrotor unmanned aerial vehicle}
\label{sec:example-quadrotor}

\TODO{模型（平移 + 姿态，姿态在 SO(3) 上完整描述）；四级层级，Taylor 预瞄阶数为 4，姿态部分用 SO(3) 上的展开（Magnus 型），输入在 u = 0 处线性化；参考轨迹的条件；控制律推导}
